%======================================================================
\chapter{Separated sets and boundary bumping}\label{ch:bumping}
%======================================================================

This chapter collects the elementary tools used in all later chapters: separated sets, unions
of continua, nested intersections of continua, components, and the boundary bumping theorem.
Throughout, $X$ is a metric space.

\section{Separated sets}

\begin{definition}\label{def:separated}
Two sets $A,B\subseteq X$ are \emph{separated} if $\cl A\cap B=\emptyset$ and $A\cap\cl
B=\emptyset$.
\end{definition}

Separated sets are disjoint. Two disjoint closed sets are separated, and so are two disjoint open
sets ($\cl A\cap B=\emptyset$ because $X\setminus B$ is a closed set containing $A$).

\begin{lemma}\label{lem:sep-of-not-conn}
If $C\subseteq X$ is not preconnected, then $C=A\cup B$ with $A,B$ nonempty and separated.
\end{lemma}

\begin{proof}
There are open sets $U,V$ with $C\subseteq U\cup V$, $C\cap U\ne\emptyset$, $C\cap V\neq\emptyset$ and
$C\cap U\cap V=\emptyset$. Put $A=C\cap U$ and $B=C\cap V$; they are nonempty and $C=A\cup B$. If
$x\in\cl A\cap B$, then $x\in V$, and the open neighbourhood $V$ of $x\in\cl A$ contains a point $y$
of $A$; then $y\in C\cap U\cap V$, which is impossible. Symmetrically $A\cap\cl B=\emptyset$.
\end{proof}

\formal{\lean{HPC.Plane.exists_separated_of_not_isPreconnected}.}

\begin{lemma}\label{lem:side}
Let $A,B$ be separated and let $T\subseteq A\cup B$ be preconnected. Then $T\subseteq A$ or
$T\subseteq B$.
\end{lemma}

\begin{proof}
The open sets $U=X\setminus\cl B$ and $V=X\setminus\cl A$ cover $T$: a point of $T$ in $A$ is not in
$\cl B$, and a point of $T$ in $B$ is not in $\cl A$. A point of $T\cap U\cap V$ would lie in $A$ or
in $B$, hence in $\cl A$ or $\cl B$, which is impossible. Since $T$ is preconnected, $T\cap U$ or
$T\cap V$ is empty. If $T\cap V=\emptyset$, every point of $T$ lies in $\cl A$, hence not in $B$
(as $\cl A\cap B=\emptyset$), hence in $A$. The other case is symmetric.
\end{proof}

\formal{\lean{HPC.Plane.subset_side_of_isPreconnected}.}

\begin{lemma}\label{lem:union-continua}
If $A$ and $B$ are continua with $A\cap B\ne\emptyset$, then $A\cup B$ is a continuum.
\end{lemma}

\begin{proof}
$A\cup B$ is compact and nonempty. If $A\cup B$ lies in the union of two disjoint open sets
$U,V$ each meeting $A\cup B$, then $A$ lies in one of them and $B$ in one of them (each is
preconnected); they cannot be the same one, since both meet $A\cup B$; but a point of $A\cap B$
lies in both, which is impossible as $U\cap V\cap(A\cup B)=\emptyset$.
\end{proof}

\formal{\lean{HPC.Plane.isContinuum_union}.}

The next lemma is the basic tool for producing new subcontinua.

\begin{lemma}[A continuum plus one piece]\label{lem:piece}
Let $S$ be a compact metric space, $N\subseteq M$ continua in $S$, and suppose
$M\setminus N=A\cup R$ with $A,R$ separated and $A\neq\emptyset$. Then $N\cup A$ is a continuum.
\end{lemma}

\begin{proof}
\emph{$N\cup A$ is closed.} We have $\cl A\subseteq M$ since $M$ is closed. A point $x\in\cl A$ not in
$N$ lies in $M\setminus N=A\cup R$, and not in $R$ because $\cl A\cap R=\emptyset$; so $x\in A$. Thus
$\cl{N\cup A}=N\cup\cl A\subseteq N\cup A$.

\emph{$N\cup A$ is preconnected.} We use the criterion \vone{(K2)}. Let $N\cup A=F\cup G$ with $F,G$
closed and disjoint. Since $N$ is preconnected, it lies in one of them, say $N\subseteq F$ (the other
case is symmetric). Put $A_G=A\cap G$ and $E=N\cup(A\cap F)\cup R$, so that $M=A_G\cup E$. These two
sets are separated:
\begin{itemize}
\item $\cl{A_G}\subseteq G\cap\cl A$; this misses $R$ (as $\cl A\cap R=\emptyset$) and misses
$N\cup(A\cap F)\subseteq F$ (as $F\cap G=\emptyset$).
\item $\cl E\subseteq F\cup\cl R$, since $N\cup(A\cap F)\subseteq F$ and $F$ is closed; $F$ misses
$A_G\subseteq G$, and $\cl R$ misses $A_G\subseteq A$.
\end{itemize}
$E$ is nonempty (it contains $N$) and $M$ is connected, so $A_G=\emptyset$ by Lemma~\ref{lem:side}
(applied to $T=M$). Hence $A\subseteq F$ and $N\cup A\subseteq F$.

So $N\cup A$ is a nonempty closed preconnected subset of the compact space $S$, that is, a
continuum.
\end{proof}

\formal{\lean{HPC.Plane.isContinuum_union_piece} (Nadler, \emph{Continuum Theory}, 6.3). A
version for an arbitrary connected topological space is Theorem~\ref{thm:nadler63}.}

\section{Nested intersections}

A family $\mathcal C$ of sets is a \emph{chain} (for inclusion) if any two of its members are
comparable: $A\subseteq B$ or $B\subseteq A$.

\begin{lemma}\label{lem:nested}
Let $S$ be a compact metric space and $\mathcal C$ a nonempty chain of continua in $S$. Then
$\bigcap\mathcal C$ is a continuum.
\end{lemma}

\begin{proof}
Put $I=\bigcap\mathcal C$; it is closed. It is nonempty: the members of $\mathcal C$ are compact and
any finitely many of them have a nonempty intersection (the smallest of them), so by compactness
the whole family has a nonempty intersection.

Suppose $I=F\cup G$ with $F,G$ closed, disjoint and nonempty. They are compact, hence at a positive
distance $3\delta$ (\vone{Lemma~1.1}). Let $U$ and $V$ be the open $\delta$-neighbourhoods of $F$
and $G$; they are disjoint. For $C\in\mathcal C$ let $C'=C\setminus(U\cup V)$, a compact set. The
sets $C'$ form a chain, and $\bigcap_{C}C'=I\setminus(U\cup V)=\emptyset$. By compactness some finitely
many of them have empty intersection; the smallest of these, $C_0'$, is then empty. So
$C_0\subseteq U\cup V$. But $C_0\supseteq I$ meets both $U\supseteq F$ and $V\supseteq G$, contradicting
the connectedness of $C_0$.
\end{proof}

\formal{\lean{HPC.T14.isPreconnected_sInter_chain}.}

\section{Components and clopen sets}

Recall (\vtwo{Section~1.4}) that the \emph{component} of a point $p$ in a set $F$ is the largest
preconnected subset of $F$ containing $p$; it is closed when $F$ is.

\begin{lemma}\label{lem:quasi}
Let $K$ be a compact metric space and $p\in K$. The component of $p$ in $K$ is the intersection
$Q$ of all clopen subsets of $K$ containing $p$.
\end{lemma}

\begin{proof}
The component $C$ of $p$ lies in every clopen set $W\ni p$: otherwise $C\cap W$ and $C\setminus W$
would split $C$ into two disjoint nonempty relatively open sets. So $C\subseteq Q$.

For the converse it suffices to show that $Q$ is preconnected (it contains $p$, so it is then
contained in $C$). Let $Q=A\cup B$ with $A,B$ closed and disjoint, $p\in A$; we show $B=\emptyset$.
Choose disjoint open sets $U\supseteq A$, $V\supseteq B$ (for instance the sets of points nearer to
$A$ than to $B$, and conversely). The closed sets $W\setminus(U\cup V)$, $W$ clopen with $p\in W$, have
empty intersection; by compactness finitely many of them do, so $W_0=W_1\cap\dots\cap W_k$
satisfies $W_0\subseteq U\cup V$. The set $W_0\cap U$ is open, and it is closed because it equals
$W_0\setminus V$. It contains $p$, so it is one of the clopen sets defining $Q$, and $Q\subseteq U$.
Then $B\subseteq Q\cap V\subseteq U\cap V=\emptyset$.
\end{proof}

\formal{Mathlib's \lean{connectedComponent_eq_iInter_isClopen}, used inside the proofs below.}

\begin{lemma}\label{lem:clopen-split}
Let $K\subseteq X$ be compact, $p\in K$ and $O$ open, and suppose that every preconnected subset of
$K$ containing $p$ lies in $O$. Then $K=H\cup L$ with $H,L$ compact and disjoint, $p\in H$ and
$H\subseteq O$.
\end{lemma}

\begin{proof}
The component $C$ of $p$ in $K$ lies in $O$. By Lemma~\ref{lem:quasi}, applied in the compact space
$K$, the closed sets $W\setminus O$, $W$ clopen in $K$ with $p\in W$, have empty intersection; by
compactness some finite intersection $H$ of such sets $W$ satisfies $H\subseteq O$. Then $H$ is
clopen in $K$, and $H$ and $L=K\setminus H$ are compact.
\end{proof}

\formal{\lean{HPC.Plane.exists_clopen_split}.}

\begin{theorem}[Separating two closed sets in a compact set]\label{thm:closed-split}
Let $K\subseteq X$ be compact and $P,Q\subseteq K$ closed, and suppose that no preconnected subset of
$K$ meets both $P$ and $Q$. Then $K=H\cup L$ with $H,L$ closed and disjoint, $P\subseteq H$ and
$Q\subseteq L$.
\end{theorem}

\begin{proof}
For each $p\in P$, Lemma~\ref{lem:clopen-split} with $O=X\setminus Q$ gives a splitting
$K=H_p\cup L_p$ into disjoint compact sets with $p\in H_p$ and $H_p\cap Q=\emptyset$. The open sets
$X\setminus L_p$ cover the compact set $P$ (since $p\notin L_p$); choose $p_1,\dots,p_k$ with
$P\subseteq\bigcup_i(X\setminus L_{p_i})$. Put $H=H_{p_1}\cup\dots\cup H_{p_k}$ and $L=K\cap
L_{p_1}\cap\dots\cap L_{p_k}$. Both are closed, and $H\cup L=K$: a point of $K$ in no $H_{p_i}$ lies
in every $L_{p_i}$. They are disjoint because $H_{p_i}\cap L_{p_i}=\emptyset$. A point of $P$ lies
outside some $L_{p_i}$, hence in $H_{p_i}\subseteq H$. Finally $Q\cap H=\emptyset$, so $Q\subseteq L$.
\end{proof}

\formal{\lean{HPC.Plane.exists_closed_split_of_no_connected} (Moore's Theorem~44 in
\emph{Foundations of Point Set Theory}). A variant with a continuum in place of the preconnected
set, in a compact space, is \lean{HPC.Plane.closed_dichotomy}.}

\section{Boundary bumping}

We use the clopen form of boundary bumping proved in Volume~2:

\begin{quote}
\vtwo{Lemma~1.5}. \emph{Let $M$ be preconnected, $L\subseteq M$ with $L\ne M$, and let $W\subseteq L$ be
nonempty, closed, and relatively open in $L$. Then $W$ meets $\cl{M\setminus L}$.}
\end{quote}

\begin{theorem}[Boundary bumping]\label{thm:bump-closed}
Let $M\subseteq X$ be preconnected, $F\subseteq M$ compact with $F\ne M$, and $p\in F$. Then there is a
continuum $K$ with $p\in K\subseteq F$ and $K\cap\cl{M\setminus F}\ne\emptyset$. One can take for
$K$ the component of $p$ in $F$.
\end{theorem}

\begin{proof}
Let $K$ be the component of $p$ in $F$; it is closed in the compact set $F$, so it is a continuum.
Suppose $K\cap\cl{M\setminus F}=\emptyset$. The set $H=F\cap\cl{M\setminus F}$ is compact and misses
$K$. By Lemma~\ref{lem:quasi} (in the compact space $F$), the closed sets $H\cap W$, $W$ clopen in $F$
with $p\in W$, have empty intersection; by compactness there is a finite intersection $W$ of such
sets with $H\cap W=\emptyset$. This $W$ is clopen in $F$: it is closed in $X$ (closed in the compact
$F$), and $W=F\cap O$ for some open $O\subseteq X$. It is nonempty ($p\in W$). By \vtwo{Lemma~1.5}
(with $L=F$) $W$ meets $\cl{M\setminus F}$; since $W\subseteq F$, it meets $H$, a contradiction.
\end{proof}

\formal{\lean{HPC.Jones.exists_continuum_bump_compact}; for $M$ a continuum in a compact space and
$F$ closed, \lean{HPC.Plane.exists_continuum_bump_closed} (Nadler 5.4).}

\begin{corollary}[Boundary bumping, open form]\label{cor:bump-open}
Let $S$ be a compact metric space, $L$ a continuum, $O$ an open set, $p\in L\cap O$ and
$L\not\subseteq O$. Then some subcontinuum $K$ of $L$ with $p\in K\subseteq\cl{L\cap O}$ meets
$L\setminus O$.
\end{corollary}

\begin{proof}
Put $F=\cl{L\cap O}\subseteq L$. If $F=L$, take $K=L$. Otherwise Theorem~\ref{thm:bump-closed} gives a
continuum $K\ni p$ inside $F$ meeting $\cl{L\setminus F}$. Now $L\setminus F\subseteq L\setminus O$
(a point of $L\cap O$ lies in $F$), and $L\setminus O$ is closed, so $\cl{L\setminus F}\subseteq
L\setminus O$.
\end{proof}

\formal{\lean{HPC.Plane.exists_continuum_bump_open} (Nadler 5.7).}

\begin{corollary}[Boundary bumping, metric form]\label{cor:bump-ball}
Let $M$ be a continuum, $p,x\in M$, $s>0$ and $d(p,x)\ge s$. Some subcontinuum $K$ of
$M\setminus B(x,s)$ contains $p$ and meets $\bar B(x,s)$.
\end{corollary}

\begin{proof}
Apply Theorem~\ref{thm:bump-closed} to $F=M\setminus B(x,s)$, which is compact, contains $p$, and
misses $x\in M$. The closure of $M\setminus F=M\cap B(x,s)$ lies in $\bar B(x,s)$.
\end{proof}

\formal{\lean{HPC.Plane.exists_subcontinuum_meeting_closedBall}.}

\begin{example}
Let $M=[0,1]$, $F=[0,\tfrac12]\cup\set1$ and $p=0$. The component of $0$ in $F$ is $[0,\frac12]$, which
meets $\cl{M\setminus F}=[\frac12,1]$ at $\frac12$. The component of $1$ in $F$ is $\set1$, which
also meets $[\frac12,1]$: boundary bumping holds for every point of $F$.
\end{example}

\section*{Exercises}

\begin{exercise}
Show that Lemma~\ref{lem:nested} fails without compactness: find a decreasing sequence of closed
connected subsets of $\R^2$ whose intersection is not connected.
\end{exercise}

\begin{exercise}
Let $M$ be a continuum and $U\subseteq M$ a nonempty open subset of $M$ with $U\neq M$. Show that every
component of $\cl U$ meets the boundary of $U$ in $M$.
\end{exercise}

\begin{exercise}
Deduce from Theorem~\ref{thm:bump-closed} that a nondegenerate continuum contains no isolated
points, and in fact that each point of it lies in arbitrarily small nondegenerate subcontinua.
\end{exercise}

\begin{exercise}
Show that Lemma~\ref{lem:quasi} fails for the non-compact space $\set{(0,0),(0,1)}\cup\bigcup_n
\set{1/n}\times[0,1]\subseteq\R^2$ with $p=(0,0)$.
\end{exercise}
